\chapter*{Abstract}
\addcontentsline{toc}{chapter}{Abstract}

Wireless signal attenuation in indoor Long Range Wide Area Networks (\gls{lorawan}) is governed by both structural building configurations and dynamic microclimatic fluctuations, limiting the predictive reliability of standard distance-only path loss models. While integrating localized environmental covariates (such as temperature, relative humidity, carbon dioxide, barometric pressure, and PM2.5 particulate matter) into regression frameworks improves path loss estimation, existing centralized approaches necessitate the continuous uploading of raw sensor streams. This constant transmission creates severe channel congestion, threatens compliance with strict LPWAN duty-cycle regulations, and presents critical data privacy risks.

This thesis addresses these challenges by introducing a decentralized, communication-compliant \gls{fl} framework combined with on-device \gls{tinyml} inference. We deploy a compact Multilayer Perceptron (MLP) with a \textbf{Dense(9$\to$8$\to$1)} architecture containing $89$ trainable parameters on resource-constrained edge nodes. The nodes predict expected path loss ($PL$ in dB) locally from 9 multimodal covariates (linearized log-distance, wall counts, environmental metrics, and gateway \gls{snr}) and classify link quality to drive event-driven status reporting. To guarantee data privacy, raw environmental sensor measurements never leave the edge device; instead, nodes run local backpropagation on buffered data. To mitigate the severe client drift caused by non-\gls{iid} data across different rooms, we implement and evaluate three federated optimization algorithms: standard Federated Averaging (\textbf{FedAvg}), \textbf{FedProx} (utilizing proximal regularization), and \textbf{SCAFFOLD} (employing control variates). For weight updates, nodes uplink quantized \texttt{int8} parameters of exactly \textbf{89~bytes} (doubled to \textbf{178~bytes} for SCAFFOLD to transmit control variates), fitting comfortably into a single \gls{lorawan} packet at standard data rates (DR5).

We evaluate the framework using a 12-month indoor multi-sensor dataset ($2,079,534$ rows) partitioned non-\gls{iid} across 6 virtual clients by physical room location. The centralized baseline achieves an $R^2$ of \textbf{0.9205} ($\text{RMSE} = 5.32\text{~dB}$). Under challenging non-\gls{iid} conditions, standard \textbf{FedAvg} converges to an RMSE of \textbf{10.87~dB}, and \textbf{FedProx} slightly improves this to \textbf{10.70~dB} with zero additional communication overhead. Crucially, \textbf{SCAFFOLD} successfully mitigates client drift, achieving an $R^2$ of \textbf{0.8619} and an RMSE of \textbf{7.01~dB}, outperforming traditional centralized linear baselines by 1.03~dB. Daily node communication bandwidth is reduced by \textbf{11.0$\times$} for FedAvg/FedProx (from 25,920 to 2,356 bytes/day) and \textbf{10.8$\times$} for SCAFFOLD (from 25,920 to 2,408 bytes/day), corresponding to over an \textbf{84\% saving in active transmission energy} for battery-powered LPWAN nodes. Compilation confirms that the TinyML inference engine and local updates easily comply with the \gls{sram} (32~KB) and Flash (256~KB) limits of the SAMD21 microcontroller on the Arduino MKR WAN 1310. This demonstrates the feasibility of decentralized, privacy-preserving path loss modeling for massive Internet of Things deployments under strict LPWAN constraints.

\noindent\textbf{Keywords:} Federated Learning, Tiny Machine Learning, LoRaWAN, LPWAN, Path Loss Modeling, Client Drift, Non-IID Data, Edge Artificial Intelligence, Indoor Propagation, Internet of Things
